\noindent\textbf{Verbatim problem statement (from file).}
\emph{``Let $A$ be a set of $n$ points in $\R^2$ such that all pairwise distances are at least $1$ and if two distinct distances differ then they differ by at least $1$. Is the diameter of $A$ $\gg n$? Perhaps the diameter is even $\ge n-1$ for sufficiently large $n$.''}

\medskip
\noindent\textbf{Definitions and conventions.}
\begin{itemize}[leftmargin=2em]
\item $n\in\N$ with $n\ge 1$.
\item $A\subset\R^2$ with $|A|=n$.
\item For distinct $x,y\in\R^2$, the Euclidean distance is $d(x,y)=\|x-y\|_2$.
\item The \emph{diameter} of $A$ is $\diam(A)=\max\{d(x,y):x,y\in A\}$.
\item The set of pairwise distances is $D(A)=\{d(x,y):x,y\in A,\ x\ne y\}$.
\item The hypothesis ``all pairwise distances are at least $1$'' means $\min D(A)\ge 1$.
\item The hypothesis ``if two distinct distances differ then they differ by at least $1$'' means:
\[\forall r,s\in D(A)\ (r\ne s\ \Rightarrow\ |r-s|\ge 1).\]
Equivalently, $D(A)$ is a $1$-separated subset of $\R$.
\end{itemize}

\medskip
\noindent\textbf{Question.}
Does there exist an absolute constant $c>0$ such that for every such $A$ one has
\[\diam(A)\ge c\,n\ ?\]
A stronger possibility mentioned is: for all sufficiently large $n$, must $\diam(A)\ge n-1$?
